\chapter{Input reference}\label{ch:inputs}

\section{Conventions}

The code uses the following conventions \notes{§2}:
\begin{itemize}
\item dimensionless normal coordinates, $q=(a+a^\dagger)/\sqrt2$;
\item energies in cm$^{-1}$;
\item $H=\sum\omega(p^2+q^2)/2+\frac{1}{3!}\sum\Phi_{ijk}q_iq_jq_k
  +\frac{1}{4!}\sum\Phi_{ijkl}q_iq_jq_kq_l$;
\item dipole derivatives in atomic units.
\end{itemize}
Indices in the files are 1-based. Force constants can be given for any
order of the indices; each distinct set of indices must appear once.

\section{The options file \file{inp}}\label{sec:inp}

All three programs read their options from \file{inp} with the
\texttt{oct\_parser} library, the input parser of the Octopus code
(included as the submodule \file{external/oct-parser}). This has some
practical consequences.

\begin{itemize}
\item \emph{Syntax.} One option per line, \texttt{name = value}; strings in
  double quotes; \texttt{\#} starts a comment; a line ending with
  \texttt{\textbackslash} continues on the next one. Names are not case
  sensitive.
\item \emph{Defaults.} Every option has a default, and only the options
  that differ from it need to be given. The defaults are listed in
  \file{README.md} of the code.
\item \emph{\file{out.oct}.} Each program writes to \file{out.oct} every
  option it has read, with the value it used. Values not given in
  \file{inp} are marked \texttt{\# default}. This is the place to check
  what a run actually used. Each program overwrites the file, so copy it
  after each phase if you want to keep all three (as \file{run.sh} of the
  tutorials does).
\item \emph{Misspelled options.} A misspelled option name is not an error:
  the parser takes it as a variable of its own, and the program uses the
  default. Check \file{out.oct} when an option seems to have no effect.
\item \emph{Variables and arithmetic.} Any name that is not an option can
  be used as a variable, and values can be expressions with the usual
  operators and functions (\texttt{sqrt}, \texttt{exp}, \texttt{log},
  \ldots) and the constants \texttt{pi} and \texttt{e}. A variable must be
  defined before it is used. The names \texttt{x}, \texttt{y}, \texttt{z},
  \texttt{r}, \texttt{w}, \texttt{t} and \texttt{i} are reserved (case
  insensitive) and cannot be used as variables. For example:
\begin{lstlisting}
thres = 0.20             # a variable, not an option
fermithres = thres
ddthres = thres
temperature = 500        # K
kT = 0.6950348*temperature   # kT is in cm-1
\end{lstlisting}
\end{itemize}

\section{Files}

\begin{description}
\item[\file{freqsfile}] One line per mode: $\omega$, then
  $\partial\mu/\partial q$ ($x$, $y$, $z$), then two optional columns: a
  legacy scale factor (unused) and a flag. In the release, a flag of 1
  computes the fundamental of that mode and 0 disables it; a disabled mode
  also has its dipole derivative set to zero. \devonly{} The flag is
  ignored, with a warning, and the bands are chosen in \file{bandlist}.
\item[\file{phi3file}, \file{phi4file}] Lines \texttt{i j k value} and
  \texttt{i j k l value}. Entries below \opt{phi3threshold} and
  \opt{phi4threshold} (cm$^{-1}$) are dropped.
\item[\file{dmu2file}] Lines \texttt{i j dx dy dz}: second derivatives of
  the dipole.
\item[\file{dmu3file}] Optional, with \opt{dipole3}: lines
  \texttt{i j k dx dy dz}. Missing entries are taken as zero.
\item[\opt{modepopfile}] Starting states: the quantum numbers of all modes,
  then the repetition count.
\item[\opt{gradientfile}] Optional: the classical gradient at the reference
  geometry, lines \texttt{i g\_i} (written by \opt{shiftgeometry}).
\item[\opt{bandlistfile}, \opt{resonanceoverrides}] \devonly{} Edited copies of
  \file{bandlist} and \file{resonancecandidates}
  (Chapter~\ref{ch:tut-review}).
\end{description}

\section{Options}

% TODO: generate this section from a single machine-readable table, shared
% with README.md, so that the two cannot diverge.
\subsection*{Resonances}
\begin{description}
\optentry{fermithres}{0.2}{1}{Threshold on $|W/\Delta|$ for cubic (Fermi) resonances.}
\optentry{ddthres}{0.2}{1}{The same for quartic (Darling--Dennison) resonances.}
\optentry{samesignresonances}{0}{1}{Also accept terms whose detuning is a single frequency.}
\optentry{closureterms}{0}{1}{Add the closure terms $-P[S_1,PH_1]$ to the resonant couplings \notes{Eq.~163}.}
\optentry{rescriterion}{0}{1}{\devonly{} 1: decide resonances from the error of perturbation theory on energies (\opt{resenergythres}, cm$^{-1}$, following~\cite{martin1995}) and relative intensities (\opt{resintensityrel}, \opt{resdarkrel}, \opt{resdarkmin}), up to \opt{rescheckquanta} quanta.}
\optentry{resonanceoverrides}{""}{1}{\devonly{} Force resonances on or off (Chapter~\ref{ch:tut-review}).}
\end{description}
\subsection*{Geometry and dipole}
\begin{description}
\optentry{shiftgeometry}{0}{1}{Move the reference geometry to remove the linear terms; writes the shifted force field (\file{*.shifted}).}
\optentry{dipole3}{0}{1}{Third-order transformed dipole for fundamentals and hot bands.}
\end{description}
\subsection*{Bands}
\begin{description}
\optentry{bandthres}{see README}{1}{Default activation of overtone, combination and difference bands, by $|\partial^2\tilde\mu|$.}
\optentry{fundthres}{1e-6}{1}{\devonly{} Default activation of fundamentals, by $|\partial\mu|$. A dipole derivative below it is numerical noise of a symmetry-forbidden mode and is set to zero when \file{freqsfile} is read.}
\optentry{bandlistfile}{""}{1}{\devonly{} Edited band list. It is read at the start of phase 1: an unknown choice stops the run at once, and lines that match no listed band give a warning.}
\end{description}
\subsection*{Checks}
\begin{description}
\optentry{dunhamcheck}{1}{1}{Stop if a deperturbed level falls at or below a lower one (Chapter~\ref{ch:trouble}).}
\optentry{dunhamfix}{0}{1}{Make the near Fermi resonances responsible explicit.}
\optentry{residualcheck}{0}{2b}{Estimate the truncation error of the polyads.}
\end{description}
\subsection*{Polyads and diagonalisation}
\begin{description}
\optentry{hsizescalepar}{see README}{2a}{Scale of the polyad growth.}
\optentry{maxtothops}{see README}{2a}{Maximum number of resonant steps.}
\optentry{costmodel}{0}{2a}{\devonly{} Cost of one resonant step when a polyad grows: 0, $1/\max(1.5,\ \opt{hsizescalepar}\,|W/\Delta E|)$; 1, $1/(\opt{costscale}\,|W/\Delta E|+1)$. Both at least $1/\opt{maxtothops}$.}
\optentry{costscale}{6}{2a}{\devonly{} Scale for \opt{costmodel}~$=1$, calibrated on toy Fermi and Wannier--Stark ladders (\file{tools/costmodel\_toy.py}).}
\optentry{minhopcoupling}{0.5}{2a}{\devonly{} A polyad does not grow through a coupling $|W|$ below this value (cm$^{-1}$); the coupling stays in the matrix when both states are in the polyad. Set it to 0 to switch it off and get the polyads of earlier versions.}
\optentry{eigensolver}{1}{2b}{Diagonalisation library: 1 ScaLAPACK \texttt{pdsyevr}, 2 ScaLAPACK \texttt{pdsyevx}, 0 SLEPc, 3 ELPA (the fastest in the tests of Chapter~\ref{ch:tut-pyrene}). With 1 and 3, the small polyads (at most 300 states per row or column of the process grid beyond the first) go to \texttt{pdsyevx}.}
\optentry{elpasolver}{1}{2b}{With \opt{eigensolver}~=~3: ELPA's one-stage (1) or two-stage (2) solver.}
\optentry{elpablock}{64}{2b}{With \opt{eigensolver}~=~3: block size of the distributed matrices. ELPA's two-stage solver gives wrong eigenvalues, without any error, with large blocks (it did with 300); keep it at 64 or below. Its GPU kernel needs a power of 2. See Section~\ref{sec:elpa}.}
\optentry{spectral}{0}{2b}{Kernel polynomial (spectral) method for large polyads; resolution \opt{spectraldelta}.}
\optentry{groundtransitions}{0}{2b}{\devonly{} With \opt{computeground}~$=1$: in the ground-state step, which adds nothing to the spectra, diagonalise all its polyads and compute the lines of the ground state, which \texttt{groundlines} (\opt{verbosity} bit 4) writes. With 0 only the polyad of the ground state is diagonalised, for the reference energy, and \texttt{groundlines} gives a warning instead of the lines. It applies to every \opt{eigensolver} and to \opt{spectral}. It does not change the spectra. Its older name, \opt{spectraleigenground}, is still read.}
\optentry{spectralexactmax}{0}{2b}{\devonly{} With \opt{spectral}~$=1$: a starting polyad of at most this many states is diagonalised, and its eigenstates replace the energy windows of the starting state (much faster for starting states in small polyads, such as the ground state). The destination polyads are still treated spectrally. Changes the spectra within the resolution; 0 switches it off.}
\end{description}
\subsection*{Histograms}
\begin{description}
\optentry{wnumhistmin, wnumhistmax, wnumhiststep}{}{2b}{Wavenumber grid (cm$^{-1}$).}
\optentry{Ehistmin, Ehistmax, Ehiststep}{}{2b}{Bins of the energy of the starting state.}
\end{description}
\subsection*{Log output}
\begin{description}
\optentry{verbosity}{0}{all}{What the logs (\file{proclog\_n*}) contain. It is a sum of the bits of Table~\ref{tab:verbosity}, which can be given by name: for example \texttt{verbosity = timings + components} instead of \texttt{verbosity = 10}. Write each name once: a name written twice adds its bit twice, which gives another bit. \opt{verbosity} does not change the spectra. In the released version bit 4 also changes what phase~2a builds: without it, phase~2a builds no polyads of the band destinations, and the spectra are empty, so keep bit 4 set there. \devonly{} In the development version \opt{verbosity} never changes what is computed; the lines of the ground state need \opt{groundtransitions}~$=1$.}
\end{description}
\begin{table}[ht]
\centering
\caption{Bits of \opt{verbosity}.}\label{tab:verbosity}
\begin{tabular}{rlp{0.62\textwidth}}
\hline
bit & name & what is written \\
\hline
1 & \texttt{polyadlevels} & For each starting state, the basis states of its polyads, with the target states marked by \texttt{*}, and their levels. With \opt{spectral}~$=1$, a polyad that is not diagonalised says so and lists only its basis states; a starting polyad diagonalised with \opt{spectralexactmax} lists its eigenstates and their weight on the starting state. \\
2 & \texttt{timings} & Timings and polyad sizes. With \opt{spectral}~$=1$, the summary lines \texttt{spectral:}, \texttt{spectral border:} and \texttt{spectral timing:} for each starting state. \\
4 & \texttt{groundlines} & The lines of the ground state, with their assignment. In the released version this bit also makes phase~2a build the polyads of the band destinations. \devonly{} The lines need \opt{groundtransitions}~$=1$. \\
8 & \texttt{components} & The main components of the eigenstates. \\
16 & \texttt{transitions} & Every transition. \\
32 & \texttt{spectraldiag} & \devonly{} Internals of the spectral method: one line per energy window of the starting state, and one per destination polyad, with its size, Chebyshev parameters and wall time. \\
\hline
\end{tabular}
\end{table}
